\subjudul{Mengatur siapa yang boleh melihat dan mengubah data di dalam objek.}
\chapter{Visibility dan Encapsulation}

\begin{dibabini}
  \item Visibility property
  \item Encapsulation
  \item Public, protected, private
  \item Name mangling
  \item Validasi lewat property
\end{dibabini}

Di akhir Bab 1 kita membuat class \texttt{Rekening} yang menolak penarikan jika saldo tidak mencukupi. Aturan itu dijaga oleh method \texttt{tarik()}. Namun ada celah besar: siapa pun tetap bisa menulis \texttt{rek.saldo = -999} dari luar class, dan aturan yang sudah kita buat dengan susah payah terlewati begitu saja.

Bab ini membahas cara menutup celah tersebut. Kita akan mengenal \emph{visibility}, yaitu tingkat keterlihatan attribute dan method, lalu menerapkan prinsip \emph{encapsulation} agar data di dalam object hanya bisa diubah melalui jalur yang sudah kita kendalikan. Di akhir bab, kita memakai fitur \emph{property} milik Python untuk memvalidasi data dengan sintaks yang tetap ringkas.

% =====================================================================
\section{Visibility Property}
% =====================================================================

\emph{Visibility} (keterlihatan) menentukan \textbf{dari mana} sebuah attribute atau method boleh diakses. Ada tiga ``wilayah'' yang biasanya dibedakan:

\begin{itemize}
  \item \textbf{di dalam class itu sendiri}, yaitu oleh method-method milik class;
  \item \textbf{di dalam class turunan} (\emph{subclass}), yang akan dibahas di Bab 4;
  \item \textbf{di luar class}, yaitu kode lain yang memakai object tersebut.
\end{itemize}

\begin{figure}[htbp]
\centering
\begin{tikzpicture}[font=\sffamily\small]
  \fill[merahmuda!35] (0,0) rectangle (12,5.6);
  \draw[merah,line width=0.8pt] (0,0) rectangle (12,5.6);
  \node[anchor=north west,text=merah,font=\sffamily\small\bfseries] at (0.2,5.45) {Kode di luar class};
  \node[anchor=north east,text=gelap] at (11.8,5.45) {boleh memakai: \texttt{nama}};
  \fill[merahmuda] (0.7,0.45) rectangle (11.3,4.6);
  \draw[merah,line width=0.8pt] (0.7,0.45) rectangle (11.3,4.6);
  \node[anchor=north west,text=merah,font=\sffamily\small\bfseries] at (0.9,4.45) {Class turunan};
  \node[anchor=north east,text=gelap] at (11.1,4.45) {boleh memakai: \texttt{nama}, \texttt{\_nama}};
  \fill[marun] (1.4,0.9) rectangle (10.6,3.55);
  \node[anchor=north west,text=emasmuda,font=\sffamily\small\bfseries] at (1.6,3.4) {Class itu sendiri};
  \node[text=white] at (6,2.0) {boleh memakai: \texttt{nama}, \texttt{\_nama}, \texttt{\_\_nama}};
\end{tikzpicture}
\caption{Tiga wilayah akses. Semakin ke dalam, semakin banyak nama yang boleh dipakai.}
\end{figure}

Bahasa seperti Java dan C++ memakai kata kunci khusus (\texttt{public}, \texttt{protected}, \texttt{private}) dan kompilatornya benar-benar menolak akses yang melanggar. Python mengambil pendekatan berbeda. Python tidak punya kata kunci visibility; yang ada adalah \textbf{konvensi penamaan} berupa garis bawah di depan nama.

\begin{table}[htbp]
\caption{Visibility di Java/C++ dan di Python.}
\begin{tabularx}{\linewidth}{P{26mm}P{40mm}L}
\kepala \kp{Tingkat} & \kp{Java/C++} & \kp{Python} \\
Public & kata kunci \texttt{public} & nama biasa: \texttt{saldo} \\
Protected & kata kunci \texttt{protected} & satu garis bawah: \texttt{\_saldo} \\
Private & kata kunci \texttt{private} & dua garis bawah: \texttt{\_\_saldo} \\
Penegakan & dipaksa oleh kompilator & sebagian besar berupa kesepakatan antarprogrammer \\
\end{tabularx}
\end{table}

\begin{info}[``We Are All Consenting Adults'']
Komunitas Python punya semboyan ``kita semua orang dewasa yang saling sepakat''. Artinya, Python memercayai programmer untuk menghormati tanda garis bawah, alih-alih memasang pagar yang tidak bisa ditembus. Tanda \texttt{\_} dan \texttt{\_\_} adalah pesan: \emph{``ini urusan internal, jangan disentuh dari luar''}.
\end{info}

% =====================================================================
\section{Konsep Encapsulation}
% =====================================================================

\emph{Encapsulation} (enkapsulasi) adalah salah satu pilar utama OOP. Konsep ini memiliki dua sisi:

\begin{enumerate}
  \item \textbf{Pembungkusan}: data (attribute) dan perilaku yang mengolahnya (method) disatukan di dalam satu class.
  \item \textbf{Penyembunyian informasi} (\emph{information hiding}): detail internal class disembunyikan, dan dunia luar hanya berinteraksi melalui antarmuka (\emph{interface}) yang disediakan, yaitu method-method publik.
\end{enumerate}

Bayangkan mesin ATM. Kita bisa menarik uang, mengecek saldo, dan mentransfer melalui tombol-tombol yang tersedia. Kita tidak bisa membuka mesin lalu mengambil uang langsung dari brankasnya. Tombol-tombol itu adalah antarmuka publik; brankas dan mekanisme di dalamnya adalah detail internal yang disembunyikan. Setiap transaksi wajib melewati tombol, sehingga aturan seperti ``saldo harus cukup'' selalu diperiksa.

\begin{figure}[htbp]
\centering
\begin{tikzpicture}[font=\sffamily\small,
    metode/.style={rectangle,fill=merah,text=white,font=\ttfamily\small\bfseries,minimum width=22mm,minimum height=8mm}]
  \fill[merahmuda!45] (0,0) rectangle (8,4.6);
  \draw[merah,line width=1pt] (0,0) rectangle (8,4.6);
  \node[anchor=south east,text=merah,font=\ttfamily\small\bfseries] at (7.85,0.15) {object rek};
  \node[rectangle,fill=marun,text=white,align=center,minimum width=38mm,minimum height=16mm] (data) at (5,2.1)
    {\textbf{data private}\\[2pt]\texttt{\_\_saldo}\\\texttt{\_\_riwayat}};
  \node[metode] (m1) at (0,3.5) {setor()};
  \node[metode] (m2) at (0,2.1) {tarik()};
  \node[metode] (m3) at (0,0.7) {saldo};
  \node[rectangle,draw=marun,line width=0.8pt,fill=white,align=center,minimum width=24mm,minimum height=12mm] (luar) at (-4.4,2.1)
    {kode di\\luar class};
  \foreach \m in {m1,m2,m3} {
    \draw[-{Latex},line width=1pt,emas] (luar.east) -- (\m.west);
    \draw[-{Latex},abu,dashed] (\m.east) -- (data.west);
  }
  \draw[-{Latex},line width=1pt,merah,dashed] (luar.north) |- node[pos=0.75,above,text=merah,font=\sffamily\footnotesize]{akses langsung ke data} (5,5.5) -- (data.north);
  \draw[merah,line width=2.2pt] (4.72,4.32) -- (5.28,4.88) (4.72,4.88) -- (5.28,4.32);
\end{tikzpicture}
\caption{Encapsulation: kode luar hanya berinteraksi lewat antarmuka publik (merah) yang memeriksa aturan; akses langsung ke data private ditolak.}
\end{figure}

\subsection{Implementasi Encapsulation}

\begin{kode}[Rekening tanpa encapsulation]
class Rekening:
    def __init__(self, pemilik, saldo=0):
        self.pemilik = pemilik
        self.saldo = saldo

    def tarik(self, jumlah):
        if jumlah > self.saldo:
            print("Saldo tidak mencukupi.")
        else:
            self.saldo -= jumlah

rek = Rekening("Ani", 100000)
rek.tarik(500000)      # ditolak oleh method
rek.saldo = -999       # ...tetapi saldo bisa diubah langsung
rek.saldo -= 500000    # dan aturan dilewati begitu saja
print(rek.saldo)
\end{kode}

\begin{keluaran}
Saldo tidak mencukupi.
-500999
\end{keluaran}

Method \texttt{tarik()} sudah bekerja dengan benar, tetapi tidak ada gunanya jika data bisa diubah tanpa melewatinya. Encapsulation memastikan perubahan data \emph{hanya} terjadi melalui method yang sudah memeriksa aturan.

\subsection{Manfaat Encapsulation}

\begin{itemize}
  \item \textbf{Integritas data}: keadaan object selalu valid karena setiap perubahan diperiksa.
  \item \textbf{Kemudahan perubahan}: cara penyimpanan data di dalam class boleh diganti tanpa mengubah kode lain yang memakai class tersebut, selama antarmuka publiknya tetap sama.
  \item \textbf{Kemudahan pemakaian}: pengguna class cukup mengetahui method publik, tanpa perlu memahami detail di dalamnya.
  \item \textbf{Kemudahan pelacakan kesalahan}: jika data menjadi salah, sumbernya hanya bisa berasal dari sedikit method, bukan dari seluruh program.
\end{itemize}

% =====================================================================
\section{Public, Protected, dan Private secara Konseptual}
% =====================================================================

Python mengenal tiga tingkat visibility yang dibedakan dari awalan nama.

\subsection{Public}

Nama tanpa awalan garis bawah, misalnya \texttt{username}, bersifat \emph{public}: boleh dibaca, diubah, dan dipanggil dari mana saja. Semua attribute dan method yang kita tulis di Bab 1 bersifat public.

\subsection{Protected}

Nama dengan satu garis bawah, misalnya \texttt{\_email}, bersifat \emph{protected}. Maknanya: ``attribute ini untuk keperluan internal class dan turunannya''. Python \textbf{tidak} melarang akses dari luar; tanda ini murni kesepakatan. Editor dan alat pemeriksa kode (\emph{linter}) biasanya memberi peringatan jika kita mengaksesnya dari luar class.

\subsection{Private}

Nama dengan dua garis bawah di depan, misalnya \texttt{\_\_password}, bersifat \emph{private}: hanya untuk dipakai di dalam class itu sendiri. Berbeda dengan protected, Python benar-benar mempersulit akses dari luar melalui mekanisme \emph{name mangling} (dibahas di subbab berikutnya), sehingga akses biasa akan gagal.

\begin{kode}[Tiga tingkat visibility]
class Akun:
    def __init__(self, username, email, password):
        self.username = username       # public
        self._email = email            # protected
        self.__password = password     # private

    def cek_password(self, teks):
        # di dalam class, attribute private boleh diakses
        return teks == self.__password

akun = Akun("ani24", "ani@mail.com", "rahasia123")
print(akun.username)
print(akun._email)                     # bisa, tetapi melanggar kesepakatan
print(akun.cek_password("rahasia123"))
print(akun.__password)                 # gagal
\end{kode}

\begin{keluaran}
ani24
ani@mail.com
True
Traceback (most recent call last):
  ...
AttributeError: 'Akun' object has no attribute '__password'. Did you mean: 'cek_password'?
\end{keluaran}

Aturan yang sama berlaku untuk method. Method pembantu yang hanya dipakai di dalam class sebaiknya diberi awalan garis bawah, misalnya \texttt{\_format\_tanggal()} atau \texttt{\_\_enkripsi()}.

\begin{table}[htbp]
\caption{Ringkasan tingkat visibility di Python.}
\begin{tabularx}{\linewidth}{P{22mm}P{24mm}P{22mm}P{22mm}L}
\kepala \kp{Tingkat} & \kp{Contoh nama} & \kp{Dalam class} & \kp{Dari luar} & \kp{Kapan dipakai} \\
Public & \texttt{nama} & boleh & boleh & antarmuka yang memang disediakan untuk pengguna class \\
Protected & \texttt{\_nama} & boleh & bisa, tetapi tidak dianjurkan & detail internal yang mungkin dipakai class turunan \\
Private & \texttt{\_\_nama} & boleh & gagal dengan cara biasa & detail internal yang benar-benar milik class itu saja \\
\end{tabularx}
\end{table}

\begin{tips}[Protected atau Private?]
Dalam praktik sehari-hari, banyak programmer Python lebih sering memakai satu garis bawah (\texttt{\_nama}) untuk menandai data internal, dan memakai dua garis bawah hanya bila perlu menghindari bentrok nama dengan class turunan. Di buku ini kita memakai \texttt{\_\_nama} untuk data yang wajib dilindungi agar efek name mangling terlihat jelas saat belajar.
\end{tips}

% =====================================================================
\section{Name Mangling}
% =====================================================================

Mengapa \texttt{akun.\_\_password} gagal diakses padahal attribute itu jelas dibuat di \texttt{\_\_init\_\_}? Jawabannya adalah \emph{name mangling}: setiap nama di dalam class yang diawali dua garis bawah (dan tidak diakhiri dua garis bawah) otomatis diubah Python menjadi \texttt{\_NamaClass\_\_nama}.

\begin{kode}[Melihat hasil name mangling]
class Rekening:
    def __init__(self, saldo):
        self.__saldo = saldo

rek = Rekening(100000)
print(rek.__dict__)
print(rek._Rekening__saldo)     # nama aslinya setelah di-mangle
\end{kode}

\begin{keluaran}
{'_Rekening__saldo': 100000}
100000
\end{keluaran}

Attribute \texttt{\_\_saldo} sebenarnya tersimpan dengan nama \texttt{\_Rekening\_\_saldo}. Di dalam class, Python menerjemahkan \texttt{self.\_\_saldo} ke nama tersebut secara otomatis, sehingga method-method class tetap bisa memakainya dengan normal. Dari luar class, penerjemahan itu tidak terjadi, sehingga \texttt{rek.\_\_saldo} tidak ditemukan.

\begin{figure}[htbp]
\centering
\begin{tikzpicture}[font=\sffamily\small,
    kd/.style={rectangle,draw=marun,line width=0.8pt,fill=merahmuda,font=\ttfamily\small,minimum height=8mm,inner xsep=3mm},
    lbl/.style={font=\sffamily\footnotesize,text=abu}]
  \node[kd] (c1) at (0,1.3) {self.\_\_saldo};
  \node[lbl,anchor=south] at (c1.north) {di dalam class};
  \node[kd] (c2) at (6.2,1.3) {self.\_Rekening\_\_saldo};
  \node[kd] (d1) at (0,-1.3) {rek.\_\_saldo};
  \node[lbl,anchor=north] at (d1.south) {di luar class};
  \node[kd] (d2) at (6.2,-1.3) {rek.\_\_saldo};
  \node[rectangle,fill=marun,text=white,align=left,minimum height=34mm,text width=36mm,inner sep=3mm] (dict) at (11.9,0)
    {\textbf{isi \_\_dict\_\_}\\[4pt]\ttfamily '\_Rekening\_\_saldo':\\ 100000};
  \draw[-{Latex},line width=1pt,emas] (c1) -- node[above,font=\sffamily\footnotesize]{diterjemahkan} (c2);
  \draw[-{Latex},line width=1pt,emas] (c2.east) -- (dict.west |- c2.east);
  \draw[-{Latex},line width=1pt,abu] (d1) -- node[above,font=\sffamily\footnotesize]{tetap apa adanya} (d2);
  \draw[-{Latex},line width=1pt,merah,dashed] (d2.east) -- (dict.west |- d2.east);
  \draw[merah,line width=2pt] ($(d2.east)!0.5!(dict.west |- d2.east)+(-0.22,-0.22)$) -- ++(0.44,0.44)
                              ($(d2.east)!0.5!(dict.west |- d2.east)+(-0.22,0.22)$) -- ++(0.44,-0.44);
  \node[font=\sffamily\footnotesize,text=merah,anchor=north] at ($(d2.east)!0.5!(dict.west |- d2.east)+(0,-0.3)$) {tidak ditemukan};
\end{tikzpicture}
\caption{Name mangling hanya terjadi untuk kode yang ditulis di dalam class.}
\end{figure}

\begin{penting}[Name Mangling Bukan Fitur Keamanan]
Seperti terlihat di atas, data private masih bisa dibaca lewat \texttt{rek.\_Rekening\_\_saldo}. Name mangling hanya mencegah akses \emph{tidak sengaja}, bukan akses yang disengaja. Jangan pernah mengandalkannya untuk menyembunyikan data rahasia seperti password; data semacam itu harus dilindungi dengan cara lain, misalnya enkripsi.
\end{penting}

\subsection{Jebakan: Menulis \texttt{\_\_nama} dari Luar Class}

Karena penerjemahan hanya terjadi di dalam class, memberi nilai ke \texttt{rek.\_\_saldo} dari luar \textbf{tidak} mengubah saldo. Python justru membuat attribute baru bernama \texttt{\_\_saldo} apa adanya.

\begin{kode}[Attribute baru, bukan mengubah yang private]
rek.__saldo = 0
print(rek.__dict__)
\end{kode}

\begin{keluaran}
{'_Rekening__saldo': 100000, '__saldo': 0}
\end{keluaran}

\subsection{Tujuan Sebenarnya: Mencegah Bentrok Nama}

Name mangling dirancang untuk melindungi attribute sebuah class agar tidak tertimpa secara tidak sengaja oleh class turunannya. Contoh berikut memakai pewarisan yang baru akan dibahas tuntas di Bab 4; cukup perhatikan bahwa class \texttt{Anak} dibangun di atas class \texttt{Induk}.

\begin{kode}[Dua attribute private dengan nama sama]
class Induk:
    def __init__(self):
        self.__kode = "INDUK"

    def kode_induk(self):
        return self.__kode

class Anak(Induk):
    def __init__(self):
        super().__init__()
        self.__kode = "ANAK"

    def kode_anak(self):
        return self.__kode

a = Anak()
print(a.kode_induk(), a.kode_anak())
print(a.__dict__)
\end{kode}

\begin{keluaran}
INDUK ANAK
{'_Induk__kode': 'INDUK', '_Anak__kode': 'ANAK'}
\end{keluaran}

Karena di-mangle menjadi \texttt{\_Induk\_\_kode} dan \texttt{\_Anak\_\_kode}, kedua attribute hidup berdampingan tanpa saling menimpa. Seandainya kita memakai satu garis bawah (\texttt{\_kode}), nilai milik \texttt{Induk} akan tertimpa menjadi \texttt{"ANAK"}.

\begin{info}[Yang Tidak Di-mangle]
Nama yang diawali \emph{dan} diakhiri dua garis bawah, seperti \texttt{\_\_init\_\_} dan \texttt{\_\_str\_\_}, tidak terkena name mangling. Nama-nama ini adalah \emph{dunder method} milik Python dan memang harus bisa dipanggil dengan nama aslinya.
\end{info}

% =====================================================================
\section{Pengendalian Akses Attribute}
% =====================================================================

Setelah attribute dijadikan private, dunia luar tidak bisa lagi membacanya maupun mengubahnya secara langsung. Kita perlu menyediakan jalur resmi. Cara klasik yang juga dipakai di banyak bahasa lain adalah \emph{getter} dan \emph{setter}:

\begin{itemize}
  \item \textbf{Getter}: method untuk membaca nilai, biasanya bernama \texttt{get\_nama()}.
  \item \textbf{Setter}: method untuk mengubah nilai sekaligus memvalidasinya, biasanya bernama \texttt{set\_nama(nilai)}.
\end{itemize}

Dengan mengatur getter dan setter mana yang disediakan, kita bisa menentukan pola akses setiap attribute.

\begin{table}[htbp]
\caption{Pola pengendalian akses attribute.}
\begin{tabularx}{\linewidth}{P{34mm}P{40mm}L}
\kepala \kp{Pola} & \kp{Yang disediakan} & \kp{Contoh} \\
Baca-tulis & getter dan setter dengan validasi & nama pemilik rekening \\
Hanya-baca (\emph{read-only}) & getter saja & nomor rekening, NIM \\
Hanya-tulis (\emph{write-only}) & setter saja & password \\
Diubah lewat aksi & getter dan method aksi, tanpa setter & saldo yang hanya berubah lewat \texttt{setor()} dan \texttt{tarik()} \\
\end{tabularx}
\end{table}

\begin{kode}[Rekening dengan getter dan setter]
class Rekening:
    def __init__(self, pemilik, saldo_awal=0):
        if saldo_awal < 0:
            raise ValueError("Saldo awal tidak boleh negatif.")
        self.__pemilik = pemilik
        self.__saldo = saldo_awal
        self.__riwayat = []

    # getter
    def get_pemilik(self):
        return self.__pemilik

    def get_saldo(self):
        return self.__saldo

    def get_riwayat(self):
        return list(self.__riwayat)    # kirim salinan, bukan list aslinya

    # setter dengan validasi
    def set_pemilik(self, nama):
        if not nama.strip():
            raise ValueError("Nama pemilik tidak boleh kosong.")
        self.__pemilik = nama.strip()

    # method aksi: satu-satunya jalan untuk mengubah saldo
    def setor(self, jumlah):
        if jumlah <= 0:
            raise ValueError("Setoran harus lebih dari nol.")
        self.__saldo += jumlah
        self.__catat("setor", jumlah)

    def tarik(self, jumlah):
        if jumlah <= 0:
            raise ValueError("Penarikan harus lebih dari nol.")
        if jumlah > self.__saldo:
            raise ValueError("Saldo tidak mencukupi.")
        self.__saldo -= jumlah
        self.__catat("tarik", jumlah)

    # method private: hanya dipakai di dalam class
    def __catat(self, jenis, jumlah):
        self.__riwayat.append((jenis, jumlah))


rek = Rekening("Ani", 100000)
rek.setor(50000)
rek.tarik(30000)
rek.set_pemilik("Ani Lestari")
print(rek.get_pemilik(), rek.get_saldo())
print(rek.get_riwayat())
rek.tarik(999999)
\end{kode}

\begin{keluaran}
Ani Lestari 120000
[('setor', 50000), ('tarik', 30000)]
Traceback (most recent call last):
  ...
ValueError: Saldo tidak mencukupi.
\end{keluaran}

Perhatikan beberapa keputusan desain pada class di atas:

\begin{itemize}
  \item Tidak ada \texttt{set\_saldo()}. Saldo hanya boleh berubah lewat \texttt{setor()} dan \texttt{tarik()}, karena keduanya memeriksa aturan dan mencatat riwayat.
  \item \texttt{get\_riwayat()} mengembalikan \emph{salinan} list. Jika yang dikembalikan list aslinya, kode di luar bisa menambah atau menghapus isi riwayat dan encapsulation pun bocor.
  \item \texttt{\_\_catat()} adalah method private karena pencatatan riwayat adalah detail internal, bukan layanan untuk pengguna class.
\end{itemize}

\begin{info}[Tentang \texttt{raise ValueError}]
Pada bab ini, data yang tidak valid ditolak dengan \texttt{raise ValueError("pesan")}. Perintah ini menghentikan program dan menampilkan pesan kesalahan, cara yang jauh lebih aman daripada diam-diam menerima data yang salah atau hanya mencetak pesan. Cara menangkap kesalahan semacam ini agar program tidak berhenti dibahas di Bab 8, \emph{Exception Handling}.
\end{info}

% =====================================================================
\section{Validasi Data melalui Property}
% =====================================================================

Getter dan setter bekerja dengan baik, tetapi penulisannya panjang. Bandingkan dua baris berikut yang maksudnya sama, yaitu mengubah nama pemilik menjadi huruf kapital.

\begin{kode}[Getter/setter vs sintaks attribute]
rek.set_pemilik(rek.get_pemilik().upper())   # getter/setter: sulit dibaca
rek.pemilik = rek.pemilik.upper()            # sintaks attribute: ringkas
\end{kode}

Python menyediakan \emph{property} untuk mendapatkan keduanya sekaligus: sintaks attribute biasa di luar, tetapi tetap dijaga oleh method di dalam.

\subsection{Decorator \texttt{@property} dan \texttt{@nama.setter}}

Property dibuat dengan dua decorator. Decorator adalah penanda berawalan \texttt{@} yang ditulis di atas method; cara kerjanya akan dibahas di Bab 3. Untuk saat ini, cukup ikuti polanya.

\begin{langkah}{Membuat Property}
  \item Simpan nilai sebenarnya di attribute private, misalnya \texttt{self.\_\_ipk}.
  \item Tulis method getter bernama \texttt{ipk} dan beri decorator \texttt{@property}.
  \item Tulis method setter dengan nama yang \textbf{sama}, \texttt{ipk}, dan beri decorator \texttt{@ipk.setter}. Letakkan validasi di sini.
  \item Di \texttt{\_\_init\_\_}, isi nilai lewat property (\texttt{self.ipk = ipk}), bukan langsung ke attribute private, agar validasi juga berlaku saat object dibuat.
\end{langkah}

\begin{kode}[Property dengan validasi]
class Mahasiswa:
    def __init__(self, nama, ipk):
        self.nama = nama          # melewati setter nama
        self.ipk = ipk            # melewati setter ipk

    @property
    def nama(self):
        return self.__nama

    @nama.setter
    def nama(self, nilai):
        nilai = nilai.strip()
        if not nilai:
            raise ValueError("Nama tidak boleh kosong.")
        self.__nama = nilai.title()

    @property
    def ipk(self):
        return self.__ipk

    @ipk.setter
    def ipk(self, nilai):
        if not 0 <= nilai <= 4:
            raise ValueError("IPK harus di antara 0 dan 4.")
        self.__ipk = nilai

mhs = Mahasiswa("  ani lestari ", 3.5)
print(mhs.nama, mhs.ipk)     # dibaca seperti attribute biasa
mhs.ipk = 3.75               # diubah seperti attribute biasa
print(mhs.ipk)
mhs.ipk = 4.5                # ditolak oleh setter
\end{kode}

\begin{keluaran}
Ani Lestari 3.5
3.75
Traceback (most recent call last):
  ...
ValueError: IPK harus di antara 0 dan 4.
\end{keluaran}

Dari luar, \texttt{mhs.ipk} tampak seperti attribute biasa. Namun setiap kali kita menulis \texttt{mhs.ipk = \ldots}, Python sebenarnya memanggil method setter, sehingga nilai yang tidak valid langsung ditolak. Setter juga bisa \emph{merapikan} data, seperti menghapus spasi berlebih dan mengubah kapitalisasi nama.

\begin{figure}[htbp]
\centering
\begin{tikzpicture}[font=\sffamily\small,
    kd/.style={rectangle,draw=marun,line width=0.8pt,fill=merahmuda,font=\ttfamily\small,minimum height=9mm,inner xsep=3mm},
    proses/.style={rectangle,fill=merah,text=white,align=center,minimum width=30mm,minimum height=12mm},
    ok/.style={rectangle,draw=hijau,line width=0.8pt,fill=hijau!12,font=\ttfamily\small,minimum height=9mm,inner xsep=3mm},
    gagal/.style={rectangle,draw=jingga,line width=0.8pt,fill=jingga!12,font=\ttfamily\small,minimum height=9mm,inner xsep=3mm},
    panah/.style={-{Latex},line width=1pt,emas}]
  \node[font=\sffamily\small\bfseries,text=merah,anchor=west] at (-1.7,3.5) {Menulis nilai};
  \node[kd] (a1) at (0,1.9) {mhs.ipk = 4.5};
  \node[proses] (a2) at (4.6,1.9) {\texttt{@ipk.setter}\\{\footnotesize validasi}};
  \node[ok] (a3) at (9.6,2.8) {self.\_\_ipk = nilai};
  \node[gagal] (a4) at (9.6,1.0) {raise ValueError};
  \draw[panah] (a1) -- (a2);
  \draw[panah] (a2.east) -- node[above,sloped,font=\sffamily\footnotesize,text=hijau]{valid} (a3.west);
  \draw[panah] (a2.east) -- node[below,sloped,font=\sffamily\footnotesize,text=jingga]{tidak valid} (a4.west);
  \draw[garis,line width=0.6pt] (-1.7,-0.1) -- (11.8,-0.1);
  \node[font=\sffamily\small\bfseries,text=merah,anchor=west] at (-1.7,-0.6) {Membaca nilai};
  \node[kd] (b1) at (0,-1.9) {print(mhs.ipk)};
  \node[proses] (b2) at (4.6,-1.9) {\texttt{@property}\\{\footnotesize getter}};
  \node[ok] (b3) at (9.6,-1.9) {return self.\_\_ipk};
  \draw[panah] (b1) -- (b2);
  \draw[panah] (b2) -- (b3);
\end{tikzpicture}
\caption{Di balik sintaks attribute biasa, property selalu melewatkan penulisan ke setter dan pembacaan ke getter.}
\end{figure}

\begin{penting}[Jangan Langsung ke Attribute Private di \texttt{\_\_init\_\_}]
Jika di \texttt{\_\_init\_\_} kita menulis \texttt{self.\_\_ipk = ipk}, validasi di setter dilewati dan object bisa lahir dengan IPK 9.0. Selalu isi nilai awal lewat property: \texttt{self.ipk = ipk}.
\end{penting}

\subsection{Property Hanya-Baca dan Property Terhitung}

Property yang hanya memiliki getter (tanpa setter) bersifat \emph{read-only}. Property juga bisa \emph{terhitung} (\emph{computed}), yaitu nilainya tidak disimpan, melainkan dihitung dari attribute lain setiap kali dibaca.

\begin{kode}[Property terhitung]
class Suhu:
    NOL_MUTLAK = -273.15

    def __init__(self, celsius):
        self.celsius = celsius

    @property
    def celsius(self):
        return self.__celsius

    @celsius.setter
    def celsius(self, nilai):
        if nilai < Suhu.NOL_MUTLAK:
            raise ValueError("Suhu tidak boleh di bawah nol mutlak.")
        self.__celsius = nilai

    @property
    def fahrenheit(self):              # terhitung dan hanya-baca
        return self.__celsius * 9 / 5 + 32

    @property
    def kelvin(self):                  # terhitung dan hanya-baca
        return self.__celsius - Suhu.NOL_MUTLAK

s = Suhu(30)
print(s.fahrenheit, s.kelvin)
s.celsius = 100
print(s.fahrenheit)                    # ikut berubah otomatis
s.fahrenheit = 50                      # tidak ada setter
\end{kode}

\begin{keluaran}
86.0 303.15
212.0
Traceback (most recent call last):
  ...
AttributeError: property 'fahrenheit' of 'Suhu' object has no setter
\end{keluaran}

Pada contoh ini, nilai \texttt{fahrenheit} dan \texttt{kelvin} dihitung dari \texttt{celsius} terbaru setiap kali dibaca dan tidak dapat diubah secara langsung.

\subsection{Deleter dengan \texttt{@nama.deleter}}

Selain getter dan setter, property dapat memiliki \emph{deleter} yang dijalankan saat kita memakai perintah \texttt{del}. Deleter jarang dipakai, tetapi berguna jika penghapusan perlu disertai tindakan lain, misalnya mencatat log.

\begin{kode}[Deleter]
class Profil:
    def __init__(self, foto):
        self.__foto = foto

    @property
    def foto(self):
        return self.__foto

    @foto.deleter
    def foto(self):
        print(f"Menghapus foto {self.__foto}")
        self.__foto = "default.png"

p = Profil("ani.jpg")
del p.foto
print(p.foto)
\end{kode}

\begin{keluaran}
Menghapus foto ani.jpg
default.png
\end{keluaran}

\begin{table}[htbp]
\caption{Getter/setter klasik dibandingkan property.}
\begin{tabularx}{\linewidth}{P{28mm}LL}
\kepala \kp{Aspek} & \kp{Getter/setter klasik} & \kp{Property} \\
Membaca & \texttt{rek.get\_pemilik()} & \texttt{rek.pemilik} \\
Mengubah & \texttt{rek.set\_pemilik("Ani")} & \texttt{rek.pemilik = "Ani"} \\
Validasi & di method setter & di method \texttt{@pemilik.setter} \\
Hanya-baca & tidak menyediakan setter & tidak menulis \texttt{@nama.setter} \\
Gaya & umum di Java dan C++ & gaya yang dianjurkan di Python (\emph{Pythonic}) \\
\end{tabularx}
\end{table}

\subsection{Studi Kasus: Rekening dengan Property}

Kita gabungkan semua yang sudah dipelajari untuk menyempurnakan class \texttt{Rekening}: nomor rekening hanya-baca, nama pemilik bisa diubah dengan validasi, dan saldo hanya-baca karena perubahannya wajib lewat \texttt{setor()} dan \texttt{tarik()}.

\begin{kode}[rekening.py]
class Rekening:
    def __init__(self, nomor, pemilik, saldo_awal=0):
        if saldo_awal < 0:
            raise ValueError("Saldo awal tidak boleh negatif.")
        self.__nomor = nomor
        self.pemilik = pemilik            # lewat setter
        self.__saldo = saldo_awal

    @property
    def nomor(self):                      # hanya-baca
        return self.__nomor

    @property
    def pemilik(self):
        return self.__pemilik

    @pemilik.setter
    def pemilik(self, nama):
        nama = nama.strip()
        if len(nama) < 3:
            raise ValueError("Nama pemilik minimal 3 huruf.")
        self.__pemilik = nama.title()

    @property
    def saldo(self):                      # hanya-baca
        return self.__saldo

    def setor(self, jumlah):
        if jumlah <= 0:
            raise ValueError("Setoran harus lebih dari nol.")
        self.__saldo += jumlah

    def tarik(self, jumlah):
        if jumlah <= 0:
            raise ValueError("Penarikan harus lebih dari nol.")
        if jumlah > self.__saldo:
            raise ValueError("Saldo tidak mencukupi.")
        self.__saldo -= jumlah

    def __str__(self):
        return f"[{self.nomor}] {self.pemilik} - saldo Rp{self.saldo:,}"


rek = Rekening("KM-001", "ani lestari", 100000)
rek.setor(25000)
rek.tarik(50000)
print(rek)
rek.saldo = -999
\end{kode}

\begin{keluaran}
[KM-001] Ani Lestari - saldo Rp75,000
Traceback (most recent call last):
  ...
AttributeError: property 'saldo' of 'Rekening' object has no setter
\end{keluaran}

Celah yang kita temukan di awal bab kini tertutup. Saldo tetap mudah dibaca dengan \texttt{rek.saldo}, tetapi satu-satunya cara mengubahnya adalah melalui method yang sudah memeriksa aturan.

% =====================================================================
\section{Rangkuman}
% =====================================================================

\begin{itemize}
  \item \textbf{Visibility} menentukan dari mana attribute dan method boleh diakses. Python memakai konvensi garis bawah, bukan kata kunci.
  \item \textbf{Encapsulation} membungkus data dan perilaku dalam satu class sekaligus menyembunyikan detail internal di balik antarmuka publik.
  \item \textbf{Public} (\texttt{nama}) bebas diakses; \textbf{protected} (\texttt{\_nama}) adalah kesepakatan untuk penggunaan internal dan class turunan; \textbf{private} (\texttt{\_\_nama}) hanya untuk class itu sendiri.
  \item \textbf{Name mangling} mengubah \texttt{\_\_nama} menjadi \texttt{\_NamaClass\_\_nama}. Tujuannya mencegah bentrok nama dengan class turunan, bukan sebagai fitur keamanan.
  \item Akses attribute private dikendalikan lewat \textbf{getter}, \textbf{setter}, atau method aksi. Pola akses bisa baca-tulis, hanya-baca, hanya-tulis, atau diubah lewat aksi.
  \item \textbf{Property} (\texttt{@property}, \texttt{@nama.setter}, \texttt{@nama.deleter}) memberi sintaks attribute biasa dengan validasi di dalam method.
  \item Isi nilai awal di \texttt{\_\_init\_\_} lewat property agar validasi juga berlaku saat object dibuat.
\end{itemize}

% =====================================================================
\section{Kuis Teoritis}
% =====================================================================

\subsection{Pilihan Ganda}

Pilihlah satu jawaban yang paling tepat.

\begin{enumerate}
  \item Dalam konvensi Python, attribute \texttt{\_harga} (satu garis bawah) menandakan attribute yang bersifat \ldots
  \begin{enumerate}[label=\alph*.,itemsep=0pt]
    \item public
    \item protected
    \item private
    \item konstan
  \end{enumerate}

  \item Manakah yang \textbf{bukan} manfaat encapsulation?
  \begin{enumerate}[label=\alph*.,itemsep=0pt]
    \item Menjaga data object tetap valid
    \item Detail internal class bisa diubah tanpa mengganggu kode pemakainya
    \item Program otomatis berjalan lebih cepat
    \item Sumber kesalahan data lebih mudah dilacak
  \end{enumerate}

  \item Di dalam class \texttt{Produk}, attribute \texttt{self.\_\_stok} akan disimpan dengan nama \ldots
  \begin{enumerate}[label=\alph*.,itemsep=0pt]
    \item \texttt{\_\_stok}
    \item \texttt{\_stok}
    \item \texttt{\_\_Produk\_stok}
    \item \texttt{\_Produk\_\_stok}
  \end{enumerate}

  \item Tujuan utama name mangling adalah \ldots
  \begin{enumerate}[label=\alph*.,itemsep=0pt]
    \item mengenkripsi data rahasia
    \item mencegah bentrok nama attribute dengan class turunan
    \item mempercepat akses attribute
    \item membuat attribute menjadi konstan
  \end{enumerate}

  \item Perhatikan kode berikut.
\begin{kode}[Soal 5]
class A:
    def __init__(self):
        self.__x = 1
    def get_x(self):
        return self.__x

a = A()
a.__x = 99
print(a.get_x())
\end{kode}
  Keluaran kode tersebut adalah \ldots
  \begin{enumerate}[label=\alph*.,itemsep=0pt]
    \item \texttt{1}
    \item \texttt{99}
    \item \texttt{AttributeError}
    \item \texttt{None}
  \end{enumerate}

  \item Decorator yang dipakai untuk membuat getter pada property adalah \ldots
  \begin{enumerate}[label=\alph*.,itemsep=0pt]
    \item \texttt{@getter}
    \item \texttt{@property}
    \item \texttt{@attribute}
    \item \texttt{@nama.getter}
  \end{enumerate}

  \item Sebuah class memiliki \texttt{@property def luas(self)} tanpa setter. Apa yang terjadi saat kita menjalankan \texttt{objek.luas = 10}?
  \begin{enumerate}[label=\alph*.,itemsep=0pt]
    \item Nilai \texttt{luas} berubah menjadi 10
    \item Dibuat attribute baru bernama \texttt{luas}
    \item Terjadi \texttt{AttributeError}
    \item Perintah diabaikan tanpa pesan
  \end{enumerate}

  \item Perhatikan kode berikut.
\begin{kode}[Soal 8]
class Barang:
    def __init__(self, harga):
        self.__harga = harga

    @property
    def harga(self):
        return self.__harga

    @harga.setter
    def harga(self, nilai):
        if nilai < 0:
            raise ValueError("Harga negatif")
        self.__harga = nilai

b = Barang(-500)
print(b.harga)
\end{kode}
  Keluaran kode tersebut adalah \ldots
  \begin{enumerate}[label=\alph*.,itemsep=0pt]
    \item \texttt{ValueError: Harga negatif}
    \item \texttt{-500}
    \item \texttt{0}
    \item \texttt{AttributeError}
  \end{enumerate}

  \item Nama berikut yang \textbf{tidak} terkena name mangling adalah \ldots
  \begin{enumerate}[label=\alph*.,itemsep=0pt]
    \item \texttt{\_\_saldo}
    \item \texttt{\_\_hitung()}
    \item \texttt{\_\_str\_\_}
    \item \texttt{\_\_data\_rahasia}
  \end{enumerate}

  \item Sebuah method \texttt{get\_daftar()} mengembalikan \texttt{self.\_\_daftar} (list asli). Masalah yang mungkin muncul adalah \ldots
  \begin{enumerate}[label=\alph*.,itemsep=0pt]
    \item list tidak bisa dibaca dari luar
    \item kode di luar bisa mengubah isi list tanpa melewati validasi class
    \item Python menolak mengembalikan attribute private
    \item list otomatis dikosongkan
  \end{enumerate}
\end{enumerate}

\subsection{Uraian}

\begin{enumerate}
  \item Jelaskan dua sisi dari konsep encapsulation dan berikan satu analogi dari kehidupan sehari-hari selain mesin ATM.
  \item Jelaskan perbedaan protected dan private di Python. Mengapa keduanya disebut ``konseptual'' dan bukan pembatasan mutlak seperti di Java?
  \item Seorang programmer mengaku ``sudah mengamankan password'' dengan menyimpannya di attribute \texttt{\_\_password}. Apakah klaim itu benar? Jelaskan dengan mengacu pada name mangling.
  \item Jelaskan mengapa nilai awal di \texttt{\_\_init\_\_} sebaiknya diisi lewat property (\texttt{self.ipk = ipk}) dan bukan langsung ke attribute private (\texttt{self.\_\_ipk = ipk}).
  \item Untuk class \texttt{Mahasiswa} dengan attribute NIM, nama, dan password, tentukan pola akses yang tepat (baca-tulis, hanya-baca, hanya-tulis, atau lewat aksi) untuk masing-masing attribute beserta alasannya.
\end{enumerate}


% =====================================================================
\section{Praktik Mandiri}
% =====================================================================

Kerjakan praktik berikut secara berurutan. Simpan setiap praktik dalam file \texttt{.py} terpisah, jalankan, dan pastikan keluarannya sesuai dengan contoh.

\subsection{Praktik 1: Produk Toko}

Buat class \texttt{Produk} dengan ketentuan:

\begin{itemize}
  \item Property \texttt{kode} bersifat hanya-baca.
  \item Property \texttt{nama}: tidak boleh kosong; spasi di awal dan akhir dihapus.
  \item Property \texttt{harga}: harus lebih dari 0.
  \item Property \texttt{stok}: tidak boleh negatif.
  \item Property terhitung \texttt{nilai\_persediaan} yang mengembalikan \texttt{harga * stok}.
  \item Method \texttt{jual(jumlah)} yang mengurangi stok jika mencukupi, dan menolak dengan \texttt{ValueError} jika tidak.
  \item Semua nilai awal di \texttt{\_\_init\_\_} wajib diisi lewat property.
\end{itemize}

\begin{keluaran}[Contoh keluaran Praktik 1]
P-01 Buku Tulis: 40 x Rp5,000 = Rp200,000
Setelah menjual 15: stok 25, nilai Rp125,000
ValueError: Stok tidak mencukupi.
\end{keluaran}

\subsection{Praktik 2: Akun Pengguna}

Buat class \texttt{AkunPengguna} dengan ketentuan:

\begin{itemize}
  \item Property \texttt{username}: minimal 5 karakter dan tidak boleh mengandung spasi; disimpan dalam huruf kecil.
  \item Password bersifat \textbf{hanya-tulis}: sediakan setter untuk menggantinya (minimal 8 karakter dan harus mengandung angka), tetapi getter selalu mengembalikan \texttt{"********"}.
  \item Method \texttt{login(password)} yang mengembalikan \texttt{True} jika password cocok.
  \item Attribute private \texttt{\_\_gagal\_login} yang bertambah setiap login gagal. Setelah 3 kali gagal, akun terkunci dan \texttt{login()} selalu mengembalikan \texttt{False}.
  \item Property hanya-baca \texttt{terkunci}.
\end{itemize}

\begin{keluaran}[Contoh keluaran Praktik 2]
ani_24 ********
Login benar: True
Login salah 3x -> terkunci: True
Login benar setelah terkunci: False
\end{keluaran}

\subsection{Praktik 3: Refaktor dengan Encapsulation}

Ambil salah satu class yang Anda buat di Praktik Mandiri Bab 1 (\texttt{Buku}, \texttt{Mahasiswa}, atau hasil Praktik 3), lalu terapkan encapsulation.

\begin{langkah}{Langkah Pengerjaan}
  \item Tentukan pola akses setiap attribute (baca-tulis, hanya-baca, hanya-tulis, atau lewat aksi) dan tuliskan alasannya dalam komentar.
  \item Ubah attribute menjadi private, lalu sediakan property atau method aksi sesuai pola yang dipilih.
  \item Pindahkan semua validasi ke setter atau method aksi; gunakan \texttt{raise ValueError} untuk data tidak valid.
  \item Jika ada attribute berupa list, pastikan getter-nya mengembalikan salinan.
  \item Uji dengan mencoba mengubah data secara tidak sah dari luar class dan tunjukkan bahwa percobaan tersebut ditolak.
\end{langkah}

\begin{penting}[Kriteria Penilaian]
\begin{itemize}
  \item Ketepatan pemilihan visibility dan pola akses setiap attribute.
  \item Validasi berada di setter atau method aksi, dan berlaku sejak object dibuat.
  \item Tidak ada jalur untuk mengubah data penting tanpa validasi.
  \item Program berjalan, dan pengujian menunjukkan data tidak valid benar-benar ditolak.
\end{itemize}
\end{penting}
